% GENERATED FILE. Edit canonical lessons, then run npm run generate:books.
% canonical-source-hash: fnv1a64:932c9cd86f1eba39
% canonical-lessons: KA-C282-apaghata, KA-C282-ambyulens, KA-C282-ottada, KA-C282-dadi, KA-C282-ascarya

\chapter{Accident, Ambulance, Stress, Nurse, Surprise}
\label{ch:ka-accident-ambulance-stress-nurse-surprise}
\hlchaptermodality{282}

\begin{chapteropening}
\textbf{By the end of this chapter:} \emph{I can say an accident, an ambulance, stress, a nurse and surprise.}

Complete the last lesson of chapter 282: I can say an accident, an ambulance, stress, a nurse and surprise.
\end{chapteropening}

\section[\texorpdfstring{\kn{ಅಪಘಾತ} (apaghāta)}{ಅಪಘಾತ (apaghāta)}]{\kn{ಅಪಘಾತ} (apaghāta) — an accident}

\label{lesson:KA-C282-apaghata}



\begin{quote}
Before the new one: say the Kannada for to fail, then the Kannada for to get stuck.
\end{quote}

\subsection*{You'll want to know: \kn{ಅಪಘಾತ}}
\textbf{\kn{ಅಪಘಾತ}} — \emph{apaghāta} — \textquotedblleft{}an accident\textquotedblright{}.

\begin{cousinweb}
\textbf{\kn{ಅಪಘಾತ}} (\emph{apaghāta}) is Sanskrit \emph{apaghāta}, \textquotedblleft{}a blow, a striking away\textquotedblright{}.
\end{cousinweb}

\begin{grammarlens}[title={What you've built}]
One more word for this chapter.
\end{grammarlens}

\subsection*{Your turn}
\begin{itemize}
\raggedright
  \item \emph{Say it:} \emph{apaghāta}
  \item \emph{Say it:} \emph{apaghāta}, once more
  \item \emph{Say it:} say \emph{apaghāta}
  \item \emph{Recall:} say the Kannada for to fail, then the Kannada for to get stuck, then say \emph{apaghāta} again
\end{itemize}

\begin{tcolorbox}[breakable,colback=teal!4,colframe=teal!35!black,title={Before you move on}]
What does \kn{ಅಪಘಾತ} mean? (\textquotedblleft{}An accident\textquotedblright{}.) Say it once more.
\end{tcolorbox}
\section[\texorpdfstring{\kn{ಆಂಬ್ಯುಲೆನ್ಸ್} (āmbyulens)}{ಆಂಬ್ಯುಲೆನ್ಸ್ (āmbyulens)}]{\kn{ಆಂಬ್ಯುಲೆನ್ಸ್} (āmbyulens) — an ambulance}

\label{lesson:KA-C282-ambyulens}



\begin{quote}
Before the new one: say the Kannada for to get stuck, then the Kannada for an accident.
\end{quote}

\subsection*{You'll want to know: \kn{ಆಂಬ್ಯುಲೆನ್ಸ್}}
\textbf{\kn{ಆಂಬ್ಯುಲೆನ್ಸ್}} — \emph{āmbyulens} — \textquotedblleft{}an ambulance\textquotedblright{}.

\begin{cousinweb}
\textbf{\kn{ಆಂಬ್ಯುಲೆನ್ಸ್}} (\emph{āmbyulens}) is a loan from English \emph{ambulance}.
\end{cousinweb}

\begin{grammarlens}[title={What you've built}]
One more word for this chapter.
\end{grammarlens}

\subsection*{Your turn}
\begin{itemize}
\raggedright
  \item \emph{Say it:} \emph{āmbyulens}
  \item \emph{Say it:} \emph{āmbyulens}, once more
  \item \emph{Say it:} say \emph{āmbyulens}
  \item \emph{Recall:} say the Kannada for to get stuck, then the Kannada for an accident, then say \emph{āmbyulens} again
\end{itemize}

\begin{tcolorbox}[breakable,colback=teal!4,colframe=teal!35!black,title={Before you move on}]
What does \kn{ಆಂಬ್ಯುಲೆನ್ಸ್} mean? (\textquotedblleft{}An ambulance\textquotedblright{}.) Say it once more.
\end{tcolorbox}
\section[\texorpdfstring{\kn{ಒತ್ತಡ} (ottaḍa)}{ಒತ್ತಡ (ottaḍa)}]{\kn{ಒತ್ತಡ} (ottaḍa) — stress, pressure}

\label{lesson:KA-C282-ottada}



\begin{quote}
Before the new one: say the Kannada for an accident, then the Kannada for an ambulance.
\end{quote}

\subsection*{You'll want to know: \kn{ಒತ್ತಡ}}
\textbf{\kn{ಒತ್ತಡ}} — \emph{ottaḍa} — \textquotedblleft{}stress, pressure\textquotedblright{}.

\begin{cousinweb}
\textbf{\kn{ಒತ್ತಡ}} (\emph{ottaḍa}) is native Dravidian, one word with Telugu \emph{ottu} and Malayalam \emph{ottanam} (DEDR 1021). DEDR lists \emph{ottaḍa}, \textquotedblleft{}pressure\textquotedblright{}, beside \emph{ottu}, \textquotedblleft{}to press\textquotedblright{}.
\end{cousinweb}

\begin{grammarlens}[title={What you've built}]
One more word for this chapter.
\end{grammarlens}

\subsection*{Your turn}
\begin{itemize}
\raggedright
  \item \emph{Say it:} \emph{ottaḍa}
  \item \emph{Say it:} \emph{ottaḍa}, once more
  \item \emph{Say it:} say \emph{ottaḍa}
  \item \emph{Recall:} say the Kannada for an accident, then the Kannada for an ambulance, then say \emph{ottaḍa} again
\end{itemize}

\begin{tcolorbox}[breakable,colback=teal!4,colframe=teal!35!black,title={Before you move on}]
What does \kn{ಒತ್ತಡ} mean? (\textquotedblleft{}Stress, pressure\textquotedblright{}.) Say it once more.
\end{tcolorbox}
\section[\texorpdfstring{\kn{ದಾದಿ} (dādi)}{ದಾದಿ (dādi)}]{\kn{ದಾದಿ} (dādi) — a nurse}

\label{lesson:KA-C282-dadi}



\begin{quote}
Before the new one: say the Kannada for an ambulance, then the Kannada for stress.
\end{quote}

\subsection*{You'll want to know: \kn{ದಾದಿ}}
\textbf{\kn{ದಾದಿ}} — \emph{dādi} — \textquotedblleft{}a nurse\textquotedblright{}.

\begin{cousinweb}
Where \textbf{\kn{ದಾದಿ}} (\emph{dādi}) comes from is uncertain. It may belong with Sanskrit \emph{dhātrī}, \textquotedblleft{}nurse\textquotedblright{} (CDIAL 6774), but no source records the Kannada word.
\end{cousinweb}

\begin{grammarlens}[title={What you've built}]
One more word for this chapter.
\end{grammarlens}

\subsection*{Your turn}
\begin{itemize}
\raggedright
  \item \emph{Say it:} \emph{dādi}
  \item \emph{Say it:} \emph{dādi}, once more
  \item \emph{Say it:} say \emph{dādi}
  \item \emph{Recall:} say the Kannada for an ambulance, then the Kannada for stress, then say \emph{dādi} again
\end{itemize}

\begin{tcolorbox}[breakable,colback=teal!4,colframe=teal!35!black,title={Before you move on}]
What does \kn{ದಾದಿ} mean? (\textquotedblleft{}A nurse\textquotedblright{}.) Say it once more.
\end{tcolorbox}
\section[\texorpdfstring{\kn{ಆಶ್ಚರ್ಯ} (āścarya)}{ಆಶ್ಚರ್ಯ (āścarya)}]{\kn{ಆಶ್ಚರ್ಯ} (āścarya) — surprise, wonder}

\label{lesson:KA-C282-ascarya}



\begin{quote}
Before the new one: say the Kannada for stress, then the Kannada for a nurse.
\end{quote}

\subsection*{You'll want to know: \kn{ಆಶ್ಚರ್ಯ}}
\textbf{\kn{ಆಶ್ಚರ್ಯ}} — \emph{āścarya} — \textquotedblleft{}surprise, wonder\textquotedblright{}.

\begin{cousinweb}
\textbf{\kn{ಆಶ್ಚರ್ಯ}} (\emph{āścarya}) is Sanskrit \emph{āścarya}, \textquotedblleft{}wonderful\textquotedblright{}, taken over whole as a learned loan, a tatsama (CDIAL 1464).
\end{cousinweb}

\begin{grammarlens}[title={What you've built}]
One more word for this chapter.
\end{grammarlens}

\subsection*{Your turn}
\begin{itemize}
\raggedright
  \item \emph{Say it:} \emph{āścarya}
  \item \emph{Say it:} \emph{āścarya}, once more
  \item \emph{Say it:} say \emph{āścarya}
  \item \emph{Recall:} say the Kannada for stress, then the Kannada for a nurse, then say \emph{āścarya} again
\end{itemize}

\begin{tcolorbox}[breakable,colback=teal!4,colframe=teal!35!black,title={Before you move on}]
What does \kn{ಆಶ್ಚರ್ಯ} mean? (\textquotedblleft{}Surprise, wonder\textquotedblright{}.) Say it once more.
\end{tcolorbox}
